\section{Positive lifting and hard-core particle sectors}
\label{sec:hastings05_tree_sector_gap}

The following estimates enter the many-body part of the tree example
concerning the unrestricted metric-lattice summability assertion in equation
(oibu) of \cite[lines 1690--1694]{Hastings2006Solving}. They apply to an
arbitrary coefficient matrix satisfying the stated one-particle lower bound.
The specific tree incidence estimate and its complete ground-space
classification are discussed in \cite{gap:hastings05_tree_basis_counterexample}.

\begin{theorem}[Positive quadratic lifting]
    \label{thm:hastings05_positive_quadratic_lift}
    \lean{TNLean.Lattice.TreeBasisCounterexample.quadraticLift,
        TNLean.Lattice.TreeBasisCounterexample.modeOperator,
        TNLean.Lattice.TreeBasisCounterexample.quadraticLift_outer,
        TNLean.Lattice.TreeBasisCounterexample.quadraticLift_gram,
        TNLean.Lattice.TreeBasisCounterexample.posSemidef_quadraticLift,
        TNLean.Lattice.TreeBasisCounterexample.quadraticLift_one,
        TNLean.Lattice.TreeBasisCounterexample.quadraticLift_gap}
    \leanok
    For a finite family of operators $s_v$, define
    $\mathcal L(A)=\sum_{u,v}A_{uv}s_u^*s_v$ and
    $B_w=\sum_vw_vs_v$. This linear map preserves positivity. More
    precisely, for every rectangular matrix $D$,
    \begin{align*}
        \mathcal L(D^*D)=\sum_k B_{D_k}^*B_{D_k}.
    \end{align*}
    Also $\mathcal L(1)=\sum_vs_v^*s_v$ and
    $\mathcal L(\overline w w^{\mathsf T})=B_w^*B_w$.
    Consequently
    \begin{align*}
        A\ge\delta(1-\overline w w^{\mathsf T})
        \quad\Longrightarrow\quad
        \mathcal L(A)\ge\delta\Bigl(\sum_vs_v^*s_v-B_w^*B_w\Bigr).
    \end{align*}
\end{theorem}
\begin{proof}\leanok
    Expand the matrix products and the finite sums. Every positive
    coefficient matrix has a factorization $D^*D$, and each lifted
    row contributes the positive operator $B_{D_k}^*B_{D_k}$.
\end{proof}

\begin{theorem}[Actual hard-core particle number]
    \label{thm:hastings05_hard_core_number}
    \lean{TNLean.Lattice.TreeBasisCounterexample.qubitLowering,
        TNLean.Lattice.TreeBasisCounterexample.qubitLowering_number,
        TNLean.Lattice.TreeBasisCounterexample.lowering,
        TNLean.Lattice.TreeBasisCounterexample.numberOperator,
        TNLean.Lattice.TreeBasisCounterexample.occupation,
        TNLean.Lattice.TreeBasisCounterexample.lowering_number,
        TNLean.Lattice.TreeBasisCounterexample.numberOperator_eq_diagonal,
        TNLean.Lattice.TreeBasisCounterexample.isStarProjection_lowering_number,
        TNLean.Lattice.TreeBasisCounterexample.norm_lowering_le_one,
        TNLean.Lattice.TreeBasisCounterexample.norm_lowering_apply_le,
        TNLean.Lattice.TreeBasisCounterexample.IsParticleSector,
        TNLean.Lattice.TreeBasisCounterexample.re_inner_numberOperator,
        TNLean.Lattice.TreeBasisCounterexample.sum_norm_lowering_sq_of_sector}
    \leanok
    Let $s_v$ be the qubit lowering matrix $|0\rangle\langle1|$ at site
    $v$ on the physical configuration space $\{0,1\}^V$.
    The number operator $\mathcal N=\sum_vs_v^*s_v$ is exactly diagonal
    with value $\sum_v\sigma_v$ on configuration $\sigma$.
    Every $s_v$ is a contraction. If $\mathcal N\phi=m\phi$, then
    \begin{align*}
        \sum_v\|s_v\phi\|^2=m\|\phi\|^2.
    \end{align*}
\end{theorem}
\begin{proof}\leanok
    For one qubit, $s^*s=\operatorname{diag}(0,1)$ is an orthogonal
    projection. The single-site embedding preserves products and adjoints,
    so summing gives the asserted physical diagonal. Taking its quadratic
    form gives the sector identity. The projection identity also implies
    $\|s_v\|\le1$.
\end{proof}

\begin{theorem}[The multi-particle sector energy bound]
    \label{thm:hastings05_hard_core_sector_gap}
    \lean{TNLean.Lattice.TreeBasisCounterexample.norm_mode_apply_le,
        TNLean.Lattice.TreeBasisCounterexample.norm_mode_apply_sq_le,
        TNLean.Lattice.TreeBasisCounterexample.norm_hardCoreMode_apply_sq_le,
        TNLean.Lattice.TreeBasisCounterexample.hardCore_sector_energy_ge}
    \leanok
    \uses{thm:hastings05_positive_quadratic_lift,
        thm:hastings05_hard_core_number}
    Let $\sum_v|w_v|^2=1$ and choose a site $r$ with
    $\sum_{v\ne r}|w_v|^2\le1/64$. For the actual qubit lowerings,
    every $m$-particle vector $\phi$, with $m\ge2$, satisfies
    \begin{align*}
        \|B_w\phi\|&\le\tfrac78\sqrt m\,\|\phi\|,\\
        \|B_w\phi\|^2&\le\tfrac{49}{64}m\|\phi\|^2.
    \end{align*}
    If the coefficient matrix satisfies
    $A\ge\tfrac{49}{64}(1-\overline w w^{\mathsf T})$, then
    \begin{align*}
        \operatorname{Re}\langle\phi,\mathcal L(A)\phi\rangle
           \ge\tfrac{735}{2048}\|\phi\|^2.
    \end{align*}
    The assumed inequality is solely on the one-particle coefficient
    matrix; the multi-particle energy bound is derived.
\end{theorem}
\begin{proof}\leanok
    Separate the site $r$ and apply Cauchy--Schwarz to the tail:
    \begin{align*}
        \|B_w\phi\|\le\|\phi\|+
          \tfrac18\Bigl(\sum_{v\ne r}\|s_v\phi\|^2\Bigr)^{1/2}
        \le(1+\sqrt m/8)\|\phi\|.
    \end{align*}
    Since $m\ge2$, $\sqrt m\ge4/3$, giving the stated contraction.
    Positive lifting then gives energy at least
    $(49/64)(15m/64)\|\phi\|^2$, whose minimum for $m\ge2$ is
    $735\|\phi\|^2/2048$.
\end{proof}
